% ---------------------------------------------------------------------
\section{Introduction}
% ---------------------------------------------------------------------

The 2024 US presidential election generated an unprecedented volume of digital discourse. Decentralized prediction markets such as Polymarket simultaneously aggregated the beliefs of financially motivated participants into a single daily probability estimate for each candidate's chances of winning. The central question of this project is whether publicly available signals from social media, news coverage, polling data, financial markets, and search behavior can be used to anticipate the daily movements of such a market price.

Prediction markets are theorized to be highly efficient aggregators of dispersed information, rapidly incorporating new signals into prices \citep{manski2006}. If social media and news data carry information that reaches these markets with a measurable delay, then the text mining, sentiment analysis, and network analysis techniques covered in this course may yield directional forecasting value. The project applies the full analytical pipeline from Lecture~1 through Lecture~6 to this real-world forecasting problem, covering data collection, preprocessing, word analysis, sentiment analysis, network analysis, and predictive modeling. Sections~2 through~7 follow this structure in order; Section~8 discusses key findings and limitations.
